% GENERATED FILE. Edit canonical lessons, then run npm run generate:books.
% canonical-source-hash: fnv1a64:abd164e385ffe47a
% canonical-lessons: MR-C336-leni, MR-C336-guha, MR-C336-gupta, MR-C336-bhuyar, MR-C336-khajina

\chapter{Caves and Hidden Treasure}
\label{ch:mr-caves-and-hidden-treasure}
\hlchaptermodality{336}

\begin{chapteropening}
\textbf{By the end of this chapter:} \emph{I can tell of caves, secret passages and treasure.}

Complete the last lesson of chapter 336: I can tell of caves, secret passages and treasure.
\end{chapteropening}

\section[\texorpdfstring{\mr{लेणी} (leṇī)}{लेणी (leṇī)}]{\mr{लेणी} (leṇī) — rock-cut caves}

\label{lesson:MR-C336-leni}



\begin{quote}
Before the new one: say the Marathi for to dig, then the Marathi for to carve.
\end{quote}

\subsection*{You'll want to know: \mr{लेणी}}
\textbf{\mr{लेणी}} — \emph{leṇī} — \textquotedblleft{}rock-cut caves\textquotedblright{}.

\begin{quote}
\mr{डोंगरात} \mr{कोरलेली} \mr{प्राचीन} \mr{लेणी} \mr{आम्ही} \mr{पाहिली}.
\end{quote}

\begin{quote}
\emph{ḍoṅgrāt korlelī prācīn leṇī āmhī pāhilī.}
\end{quote}

\begin{quote}
We saw the ancient caves carved into the hill.
\end{quote}

Neuter.

A neuter plural, the usual form: \mr{ही} \mr{लेणी}, these caves; one cave is \mr{लेणं}.

\begin{cousinweb}
\textbf{\mr{लेणी}} (\emph{leṇī}) is usually traced to Sanskrit \emph{layana}, 'a place of rest, a cell', which Pali has as \emph{lēṇa}, 'a rock cell used by ascetics', and Prakrit as \emph{lēṇa}, 'a cave' (CDIAL 10963); CDIAL gives no Marathi form, so the path is not settled.
\end{cousinweb}

\begin{grammarlens}[title={What you've built}]
One more word for telling what happened.
\end{grammarlens}

\subsection*{Your turn}
\begin{itemize}
\raggedright
  \item \emph{Say it:} \emph{leṇī}
  \item \emph{Say it:} \emph{leṇī}, once more
  \item \emph{Say it:} \emph{ḍoṅgrāt korlelī prācīn leṇī āmhī pāhilī.} — We saw the ancient caves carved into the hill.
  \item \emph{Recall:} say the Marathi for to dig, then the Marathi for to carve, then say \emph{leṇī} again
\end{itemize}

\begin{tcolorbox}[breakable,colback=teal!4,colframe=teal!35!black,title={Before you move on}]
What does \mr{लेणी} mean? (\textquotedblleft{}Rock-cut caves\textquotedblright{}.) Say it once more.
\end{tcolorbox}
\section[\texorpdfstring{\mr{गुहा} (guhā)}{गुहा (guhā)}]{\mr{गुहा} (guhā) — a cave}

\label{lesson:MR-C336-guha}



\begin{quote}
Before the new one: say the Marathi for to carve, then the Marathi for rock-cut caves.
\end{quote}

\subsection*{You'll want to know: \mr{गुहा}}
\textbf{\mr{गुहा}} — \emph{guhā} — \textquotedblleft{}a cave\textquotedblright{}.

\begin{quote}
\mr{वाघ} \mr{गुहेत} \mr{झोपला} \mr{होता}.
\end{quote}

\begin{quote}
\emph{vāgh guhet jhoplā hotā.}
\end{quote}

\begin{quote}
The tiger was asleep in the cave.
\end{quote}

Feminine.

The oblique is \mr{गुहे}-: \mr{गुहेत}, in the cave.

\begin{cousinweb}
\textbf{\mr{गुहा}} (\emph{guhā}) is a tatsama, taken from Sanskrit \emph{guhā}, 'a hiding place, a cave' (Monier-Williams).
\end{cousinweb}

\begin{grammarlens}[title={What you've built}]
One more word for telling what happened.
\end{grammarlens}

\subsection*{Your turn}
\begin{itemize}
\raggedright
  \item \emph{Say it:} \emph{guhā}
  \item \emph{Say it:} \emph{guhā}, once more
  \item \emph{Say it:} \emph{vāgh guhet jhoplā hotā.} — The tiger was asleep in the cave.
  \item \emph{Recall:} say the Marathi for to carve, then the Marathi for rock-cut caves, then say \emph{guhā} again
\end{itemize}

\begin{tcolorbox}[breakable,colback=teal!4,colframe=teal!35!black,title={Before you move on}]
What does \mr{गुहा} mean? (\textquotedblleft{}A cave\textquotedblright{}.) Say it once more.
\end{tcolorbox}
\section[\texorpdfstring{\mr{गुप्त} (gupta)}{गुप्त (gupta)}]{\mr{गुप्त} (gupta) — secret; hidden}

\label{lesson:MR-C336-gupta}



\begin{quote}
Before the new one: say the Marathi for rock-cut caves, then the Marathi for a cave.
\end{quote}

\subsection*{You'll want to know: \mr{गुप्त}}
\textbf{\mr{गुप्त}} — \emph{gupta} — \textquotedblleft{}secret; hidden\textquotedblright{}.

\begin{quote}
\mr{त्याने} \mr{गुप्त} \mr{जागी} \mr{पैसे} \mr{ठेवले}.
\end{quote}

\begin{quote}
\emph{tyāne gupta jāgī paise ṭhevle.}
\end{quote}

\begin{quote}
He kept the money in a secret place.
\end{quote}

\mr{गुप्त} keeps one form: \mr{गुप्त} \mr{मार्ग}, \mr{गुप्त} \mr{गोष्ट}.

\begin{cousinweb}
\textbf{\mr{गुप्त}} (\emph{gupta}) is a tatsama, taken from Sanskrit \emph{gupta}, 'guarded, protected, hidden' (Monier-Williams).
\end{cousinweb}

\begin{grammarlens}[title={What you've built}]
One more word for telling what happened.
\end{grammarlens}

\subsection*{Your turn}
\begin{itemize}
\raggedright
  \item \emph{Say it:} \emph{gupta}
  \item \emph{Say it:} \emph{gupta}, once more
  \item \emph{Say it:} \emph{tyāne gupta jāgī paise ṭhevle.} — He kept the money in a secret place.
  \item \emph{Recall:} say the Marathi for rock-cut caves, then the Marathi for a cave, then say \emph{gupta} again
\end{itemize}

\begin{tcolorbox}[breakable,colback=teal!4,colframe=teal!35!black,title={Before you move on}]
What does \mr{गुप्त} mean? (\textquotedblleft{}Secret; hidden\textquotedblright{}.) Say it once more.
\end{tcolorbox}
\section[\texorpdfstring{\mr{भुयार} (bhuyār)}{भुयार (bhuyār)}]{\mr{भुयार} (bhuyār) — an underground passage}

\label{lesson:MR-C336-bhuyar}



\begin{quote}
Before the new one: say the Marathi for a cave, then the Marathi for secret.
\end{quote}

\subsection*{You'll want to know: \mr{भुयार}}
\textbf{\mr{भुयार}} — \emph{bhuyār} — \textquotedblleft{}an underground passage\textquotedblright{}.

\begin{quote}
\mr{किल्ल्यातून} \mr{नदीपर्यंत} \mr{एक} \mr{गुप्त} \mr{भुयार} \mr{होतं}.
\end{quote}

\begin{quote}
\emph{killyātūn nadīparyant ek gupta bhuyār hotaṁ.}
\end{quote}

\begin{quote}
There was a secret passage from the fort to the river.
\end{quote}

Neuter.

\mr{भुयारी} \mr{मार्ग}, a subway.

\begin{cousinweb}
\textbf{\mr{भुयार}} (\emph{bhuyār}) is the Marathi word for a passage under the ground; its origin is not settled.
\end{cousinweb}

\begin{grammarlens}[title={What you've built}]
One more word for telling what happened.
\end{grammarlens}

\subsection*{Your turn}
\begin{itemize}
\raggedright
  \item \emph{Say it:} \emph{bhuyār}
  \item \emph{Say it:} \emph{bhuyār}, once more
  \item \emph{Say it:} \emph{killyātūn nadīparyant ek gupta bhuyār hotaṁ.} — There was a secret passage from the fort to the river.
  \item \emph{Recall:} say the Marathi for a cave, then the Marathi for secret, then say \emph{bhuyār} again
\end{itemize}

\begin{tcolorbox}[breakable,colback=teal!4,colframe=teal!35!black,title={Before you move on}]
What does \mr{भुयार} mean? (\textquotedblleft{}An underground passage\textquotedblright{}.) Say it once more.
\end{tcolorbox}
\section[\texorpdfstring{\mr{खजिना} (khajinā)}{खजिना (khajinā)}]{\mr{खजिना} (khajinā) — a treasure}

\label{lesson:MR-C336-khajina}



\begin{quote}
Before the new one: say the Marathi for secret, then the Marathi for an underground passage.
\end{quote}

\subsection*{You'll want to know: \mr{खजिना}}
\textbf{\mr{खजिना}} — \emph{khajinā} — \textquotedblleft{}a treasure\textquotedblright{}.

\begin{quote}
\mr{भुयारात} \mr{राजाचा} \mr{खजिना} \mr{लपवला} \mr{होता}.
\end{quote}

\begin{quote}
\emph{bhuyārāt rājācā khajinā lapavlā hotā.}
\end{quote}

\begin{quote}
The king's treasure had been hidden in the passage.
\end{quote}

Masculine.

Also a treasury.

\begin{cousinweb}
\textbf{\mr{खजिना}} (\emph{khajinā}) is a loan from Persian \emph{xazīna}, usually traced to Arabic \emph{khazīna}, 'a treasury'.
\end{cousinweb}

\begin{grammarlens}[title={What you've built}]
One more word for telling what happened.
\end{grammarlens}

\subsection*{Your turn}
\begin{itemize}
\raggedright
  \item \emph{Say it:} \emph{khajinā}
  \item \emph{Say it:} \emph{khajinā}, once more
  \item \emph{Say it:} \emph{bhuyārāt rājācā khajinā lapavlā hotā.} — The king's treasure had been hidden in the passage.
  \item \emph{Recall:} say the Marathi for secret, then the Marathi for an underground passage, then say \emph{khajinā} again
\end{itemize}

\begin{tcolorbox}[breakable,colback=teal!4,colframe=teal!35!black,title={Before you move on}]
What does \mr{खजिना} mean? (\textquotedblleft{}A treasure\textquotedblright{}.) Say it once more.
\end{tcolorbox}
